\documentclass[aspectratio=169, 10pt]{beamer}

% --- Paquetes Modernos ---
\usepackage[utf8]{inputenc}
\usepackage[spanish,es-tabla]{babel}
\usepackage{graphicx}
\usepackage{xcolor}
\usepackage{booktabs}
\usepackage{siunitx}
\usepackage[american]{circuitikz}
\usepackage{listings}

% --- Configuración Visual UCB ---
\usetheme{Boadilla}
\usecolortheme{seahorse}

\sisetup{
    output-decimal-marker = {,},
    per-mode = symbol
}

\lstset{
    language=Python,
    basicstyle=\ttfamily\scriptsize,
    keywordstyle=\color{blue}\bfseries,
    stringstyle=\color{teal},
    commentstyle=\color{gray}\itshape,
    backgroundcolor=\color{gray!5},
    frame=single,
    breaklines=true
}

% --- Metadatos del Documento ---
\title[Introducción a la Mecatrónica]{Modelado de Transductores Térmicos}
\subtitle{Redes de Impedancia y Efectos Parásitos en CA}
\author[M.Sc. Ing. B. Quiroga]{M.Sc. Ing. Bernardo Quiroga Turdera}
\institute[UCB Tarija]{
    Universidad Católica Boliviana ``San Pablo''\\
    Sede Tarija - Carrera de Ingeniería Mecatrónica
}
\date{10 de agosto de 2026}

\begin{document}

\begin{frame}
    \titlepage
\end{frame}

\begin{frame}{Objetivos de la Sesión 03}
    \begin{itemize}
        \item \textbf{Matemática:} Generalizar las Leyes de Kirchhoff al dominio fasorial.
        \item \textbf{Diseño (CDIO):} Modelar el puente de Wheatstone con parásitos de línea.
        \item \textbf{Software:} Desarrollar un script en Python (estructuras de datos \texttt{pandas}) para extraer el desfase de señales medidas experimentalmente.
    \end{itemize}
\end{frame}

\begin{frame}{Álgebra de Impedancias Complejas}
    La abstracción del dominio temporal al fasorial permite utilizar álgebra compleja para circuitos reactivos.
    
    \vspace{0.3cm}
    \textbf{Ley de Ohm Generalizada:}
    \begin{equation*}
        \hat{V} = Z \cdot \hat{I} \implies Z = \frac{\hat{V}}{\hat{I}} = \vert{}Z\vert{} \angle \theta = R + jX
    \end{equation*}

    \vspace{0.3cm}
    \textbf{Naturaleza de la Reactancia ($X$):}
    \begin{itemize}
        \item $X > 0 \implies$ Inductiva ($X_L = \omega L$)
        \item $X < 0 \implies$ Capacitiva ($X_C = -\frac{1}{\omega C}$)
    \end{itemize}
\end{frame}

\begin{frame}{Modelado del Sensor NTC con Efectos Parásitos}
    \begin{columns}
        \begin{column}{0.45\textwidth}
            \begin{circuitikz}[scale=0.75, transform shape]
                \draw (0,4) to[sV, l=$V_{in}$] (0,0);
                \draw (0,4) -- (2,4) -- (4.5,4);
                \draw (0,0) -- (2,0) -- (4.5,0);
                \draw (2,0) node[ground]{};
                
                \draw (2,4) to[R, l=$R_1$] (2,2) to[R, l=$R_2$] (2,0);
                \draw (4.5,4) to[R, l=$R_3$] (4.5,2) to[generic, l=$Z_{sensor}$] (4.5,0);
                
                \draw (2,2) to[short, *-o] (2.7,2) node[above]{$+$};
                \draw (4.5,2) to[short, *-o] (3.8,2) node[above]{$-$};
                \node at (3.25, 2) {$\hat{V}_{out}$};
            \end{circuitikz}
        \end{column}
        
        \begin{column}{0.5\textwidth}
            \textbf{Resistencia NTC (Steinhart-Hart):}
            \begin{equation*}
                R_{NTC}(T) = R_0 \cdot e^{\beta \left( \frac{1}{T} - \frac{1}{T_0} \right)}
            \end{equation*}
            
            \vspace{0.2cm}
            \textbf{Impedancia del Cable (15m):}
            \begin{equation*}
                Z_{sensor} = R_{cab} + j\omega L_{cab} + \left( R_{NTC} \parallel X_{C_{para}} \right)
            \end{equation*}
        \end{column}
    \end{columns}
\end{frame}

\begin{frame}[fragile]{Script Base para Laboratorio 1 (Python)}
    El procesamiento computacional calculará el desfase real de las trazas del osciloscopio:
    
\begin{lstlisting}
import numpy as np
import pandas as pd

df = pd.read_csv("lab1_medicion_30C.csv")
t, v_in, v_out = df['Time'].values, df['Channel1'].values, df['Channel2'].values

# Busqueda de indices de los picos maximos
idx_pico_in = np.argmax(v_in)
idx_pico_out = np.argmax(v_out)
delta_t = t[idx_pico_out] - t[idx_pico_in]

# Conversion de retardo temporal a grados
fase_grados = delta_t * 10000.0 * 360.0
print(f"Desfase experimental: {fase_grados:.2f} grados")
\end{lstlisting}
\end{frame}

\end{document}